\section{Conclusion}

A per-stage foreground self-calibration method was developed for a low-cost PSoC-based geophone front end. Four internal VDACs, an analog multiplexer, the acquisition delta-sigma ADC, an otherwise-idle hardware filter, and a configurable control algorithm are combined to restore the physical operating point of every analog stage. A PI controller with deadband and anti-windup provides a simple implementation, while the architecture allows other control laws and DAC realizations when a different correction range or resolution is required. Unlike an ADDER-only patch, the method also addresses offsets before and after the summing node. Unlike the continuously active servo tested during development, it is disabled during acquisition and therefore adds no servo poles or settling-time versus low-frequency-distortion compromise to the measurement path. Stored gain-dependent corrections make full recalibration infrequent, while automatic verification and recalibration preserve robustness when conditions change. The result combines analog signal conditioning in the required frequency band with a highly configurable digital correction layer, allowing inexpensive analog components to satisfy system-level operating-point requirements without adding acquisition-time feedback. \textcolor{draftred}{The final conclusion must include measured offset reduction, dynamic-range improvement, repeatability, and calibration time after completing the experiments.}

